\section{Introduction}\label{sec:intro}

Droplet impact on curved surfaces is ubiquitous in nature and industrial applications, from rain on plant leaves to spray coating and anti-icing systems. Recent experimental work by Liu et al.~\cite{liu2015symmetry} demonstrated that droplets bouncing on cylindrical superhydrophobic surfaces exhibit asymmetric spreading dynamics, with contact time reduced by up to 40\% compared to flat surfaces. This phenomenon arises from anisotropic momentum redistribution driven by the surface curvature.

Numerical simulation of such phenomena requires accurate treatment of both the two-phase flow dynamics and the wetting boundary conditions at curved solid surfaces. The lattice Boltzmann method (LBM) has emerged as a popular choice for droplet impact simulations due to its natural parallelization and ease of handling complex geometries~\cite{shan1993lattice, fakhari2017improved}. Among the various multiphase LBM approaches, the phase-field method based on the Allen-Cahn equation~\cite{lee2005lattice} offers good stability at high density ratios and sharp interface tracking.

However, a fundamental limitation exists: the Allen-Cahn sharpening term, which maintains interface sharpness, \textbf{resists} geometric wetting boundary condition corrections near curved surfaces. This resistance leads to inaccurate contact angles and suppressed asymmetric spreading in simulations. This problem has not been identified or addressed in the existing literature.

In this paper, we propose a \textbf{geometric amplification} factor $\alpha_{geo}$ that overcomes the AC sharpening resistance by over-correcting the interface gradient near the wall. We derive $\alpha_{geo}$ from the balance between AC sharpening and geometric wetting, identify three threshold conditions (contact angle, curvature ratio, and resolution) through systematic calibration, and construct the first We $\times$ $R^*$ phase diagram of asymmetric droplet spreading via Allen-Cahn LBM at 828:1 density ratio. Our results are validated against the experimental and LBM data of Liu et al.~\cite{liu2015symmetry}, demonstrating that geometric amplification is essential for Allen-Cahn LBM to accurately capture wetting physics on curved surfaces.

The remainder of this paper is organized as follows. \Cref{sec:related} reviews related work. \Cref{sec:geo_amp} describes the geometric amplification method and its calibration. \Cref{sec:experiments} presents the phase diagram and analysis. \Cref{sec:conclusion} concludes with future directions.
